% ============================================================================
%  Multizoom dwarf galaxy simulations in a 25 Mpc box
%  Draft manuscript -- AASTeX 6.3.1, ApJ format
%
%  STATUS MARKERS used throughout:
%     \todo{...}      writing not yet done
%     \pending{...}   waiting on a run or measurement still in flight
%     \check{...}     number is in hand but should be re-derived at freeze
%
%  Numbers marked \measured are taken from analysis already run; their source
%  file is given in the margin comment so every value is traceable.
% ============================================================================
\documentclass[twocolumn,linenumbers]{aastex631}

\usepackage{amsmath}
\usepackage{xcolor}

\newcommand{\todo}[1]{\textcolor{red}{\textbf{[TODO: #1]}}}
\newcommand{\pending}[1]{\textcolor{blue}{\textbf{[PENDING: #1]}}}
\newcommand{\check}[1]{\textcolor{orange}{\textbf{[CHECK: #1]}}}
\newcommand{\measured}[1]{#1}

\newcommand{\Msun}{\ensuremath{M_\odot}}
\newcommand{\Mvir}{\ensuremath{M_{200\mathrm{c}}}}
\newcommand{\Rvir}{\ensuremath{R_{200\mathrm{c}}}}
\newcommand{\Mstar}{\ensuremath{M_\star}}
\newcommand{\kpch}{\ensuremath{\mathrm{kpc}\,h^{-1}}}
\newcommand{\Mpch}{\ensuremath{\mathrm{Mpc}\,h^{-1}}}
\newcommand{\enzo}{\textsc{enzo}}
\newcommand{\music}{\textsc{music}}
\newcommand{\grackle}{\textsc{grackle}}
\newcommand{\ahf}{\textsc{ahf}}

\shorttitle{Multizoom dwarf simulations}
\shortauthors{Tumlinson et al.}

\begin{document}

\title{Many Zooms, One Box: An Efficient Method for Simulating Populations of
Dwarf Galaxies, and First Results from a 25~Mpc Campaign}

\author{J.~Tumlinson}
\affiliation{Space Telescope Science Institute, 3700 San Martin Drive,
Baltimore, MD 21218, USA}
\todo{full author list, affiliations, ORCIDs}

\begin{abstract}
\todo{Write last. Target structure, one sentence each:}
(1) Dwarf galaxies below the SAGA mass range are where quenching physics is
least constrained and where zoom simulations are most expensive per object.
(2) We present \emph{multizoom}, a method that evolves $N$ independent
Lagrangian zoom regions inside a single \enzo{} domain, motivated by the
measurement that deep levels in a single-zoom run leave
\measured{36--47\%} of rank-time idle because they hold fewer grids than MPI
ranks.
(3) Six halos evolved together cost \measured{0.40$\times$} six standalone
runs, and the cost separates cleanly into an $N$-independent root grid plus a
marginal per-halo term.
(4) We validate the method against standalone counterparts: halo masses agree
to a median ratio of \measured{0.997}, assembly histories track from $z=10$ to
$z=0$, and substructure and density profiles are indistinguishable.
(5) We apply it to a fleet of \check{28} dwarf halos spanning the
ELVES--SAGA transition and report star formation histories, scaling relations,
and the quenched fraction as a function of mass.
(6) \pending{headline science result once the gas runs complete}
\end{abstract}

\keywords{Galaxy formation (595) --- Dwarf galaxies (416) ---
Hydrodynamical simulations (767) --- N-body simulations (1083)}

% ============================================================================
\section{Introduction}
\label{sec:intro}
% ============================================================================

\todo{Full introduction. Argument skeleton:}

\begin{enumerate}
\item \textbf{The problem.} Quenching in dwarfs. Above
$\Mstar \sim 10^{9}\,\Msun$ field dwarfs are almost all star-forming; below
$\sim 10^{6}\,\Msun$ nearly all are quenched. The transition is where the
physics lives, and it is exactly where samples are thin. Cite ELVES
\citep[][and the Kado-Fong SAGA background analysis]{Carlsten2022},
SAGA \citep{Geha2017,Mao2021}, and the SAGA background papers.

\item \textbf{Why simulations struggle here.} Resolving the ISM of a
$10^{8}$--$10^{10}\,\Msun$ halo requires zoom refinement, and one zoom
produces one galaxy. Statistical statements about quenched fractions need
tens to hundreds. Cite FIRE \citep{Hopkins2018,Wetzel2016}, EDGE
\citep{Agertz2020}, Auriga, and the existing dwarf zoom suites.

\item \textbf{The cost structure nobody exploits.} A deep zoom is not
compute-bound in the usual sense --- it is \emph{granularity}-bound. We show
(\S\ref{sec:multizoom}) that the deep levels of a production dwarf zoom hold
fewer grids than MPI ranks, so most of the machine idles.

\item \textbf{This work.} We introduce multizoom, validate it, and apply it.

\item \textbf{Roadmap.} \S\ref{sec:sims} simulations; \S\ref{sec:multizoom}
the method; \S\ref{sec:validation} validation; \S\ref{sec:resolution}
resolution; \S\ref{sec:results} first results; \S\ref{sec:summary} summary.
\end{enumerate}

% ============================================================================
\section{Simulations}
\label{sec:sims}
% ============================================================================

\subsection{Code and physics}
\label{sec:code}

\enzo{} \citep{Bryan2014}, adaptive mesh refinement, with \grackle{}
\citep{Smith2017} for non-equilibrium cooling and chemistry.

\todo{Describe, with parameter values from \texttt{templates/gas-LX.enzo}:}
\begin{itemize}
\item Cosmology: $h=0.7$, $\Omega_m = 0.291$, $\Omega_\Lambda = 0.709$.
      \check{quote $\Omega_b$, $\sigma_8$, $n_s$ from the \music{} configs}
\item Star formation: H$_2$-regulated prescription. \textbf{Important
      systematic:} the prescription is bistable in this regime --- ignition
      occurs at $z\sim6$ at nref9 and $z\sim3.95$ at nref7
      (\S\ref{sec:resolution}).
\item Feedback: mechanical supernova feedback \citep[][Cassi's
      implementation]{Lochhaas2023}, \texttt{StarParticleFeedback = 64},
      with tabulated yields, stochastic SNe, and pre-SN feedback.
\item Radiation: a local stellar radiation module is compiled in but
      \emph{disabled} throughout this campaign, so that all runs remain
      mutually comparable.
\item Two-stage cooling: runs execute in two legs, switching on H$_2$
      self-shielding and the shielded \grackle{} table at $z=15$.
\end{itemize}

\subsection{Parent simulation and halo selection}
\label{sec:selection}

A $25\,\mathrm{Mpc}\,h^{-1}$ box with a $512^3$ root grid.
\todo{Describe the parent DM-only run, the $z=0$ halo catalog, and the
isolation criteria used to build the target list.}

The registry contains \check{32} halos, of which \check{28} are enabled,
selected to span the ELVES-to-SAGA transition in halo mass.
\todo{Table~\ref{tab:halos}: halo ID, \Mvir, \Rvir, isolation metric,
resolution levels run, zoom radius floor.}

\subsection{Zoom initial conditions}
\label{sec:ics}

Nested initial conditions with \music{} \citep{Hahn2011}, built level by
level: the Lagrangian region of each target is traced back from a completed
run at level $N-1$ and re-gridded at level $N$.

\paragraph{Zoom radius floor.} The comoving radius enclosing the traced
Lagrangian volume is floored at a per-halo value \texttt{rvir\_min}. This is
not cosmetic --- \S\ref{sec:contamination} shows that the floor is the single
control that determines whether a halo is contaminated at L3.

\subsection{Resolution ladder and nomenclature}
\label{sec:ladder}

\todo{Define L1/L2/L3 zoom levels and the nref cooling ceiling, with a table
of particle masses and cell sizes:}
\begin{itemize}
\item L2 fine DM particle mass \measured{$1.77\times10^{5}\,\Msun$}
\item L3 fine DM particle mass \measured{$2.62\times10^{4}\,\Msun$},
      coarse \measured{$1.34\times10^{7}\,\Msun$}
\item nref7 / nref8 / nref9 maximum refinement levels and their physical
      cell sizes at $z=0$.
\end{itemize}

\subsection{Outputs and analysis}
\label{sec:outputs}

DM-only runs write 15 redshift outputs ($z = 99, 50, 10, 9, \ldots, 0.2, 0$);
gas runs write $\sim$100.
Halos are identified with \ahf{} \citep{Knollmann2009} at $\Delta = 200$
relative to the critical density.
\todo{Describe the halocat property pipeline: Gadget export, \ahf{}, and the
streaming property engine; state the $f_\mathrm{hires}$ purity cut.}

% ============================================================================
\section{Multizoom}
\label{sec:multizoom}
% ============================================================================

\subsection{Motivation: deep zoom levels starve the machine}
\label{sec:starvation}

\todo{Present the measurement that motivates everything.} In a production
single-zoom gas run, \measured{36--47\%} of rank-time is lost purely because
deep levels hold fewer grids than there are MPI ranks. Level~7 alone accounts
for \measured{43\%} of wallclock while holding only \measured{$\sim$30 grids
on 128 ranks}. Dark matter runs do not suffer this (\measured{1.9\%}), which
is why DM multizoom gains come from sharing the root grid rather than from
occupancy.

\begin{figure}
\plotone{figures/PLACEHOLDER_rank_occupancy}
\caption{\pending{Grids per rank and idle fraction as a function of AMR level
for a production single-zoom run. This figure does not exist yet and must be
generated from \texttt{performance.out}.}}
\label{fig:starvation}
\end{figure}

\subsection{Method}
\label{sec:mz-method}

$N$ Lagrangian zoom regions are placed in a single \enzo{} domain. Each
carries its own must-refine particle set; refinement follows those particles
exactly as in a single zoom, so no new flagging machinery is required
(\texttt{CellFlaggingMethod = 4 8}).

\todo{Describe the IC construction: per-halo \music{} runs at a
\emph{shared} domain shift, then a merge step that combines the nested grids
into one hierarchy.}

\paragraph{A distinction that matters.} Multizoom ($N$ zoom \emph{regions})
is not the same mechanism as \enzo{}'s \texttt{MultiRefineRegion}
(\texttt{CellFlaggingMethod = 20}, $N$ forced-refinement \emph{boxes}). The
two are independent and compose: a multizoom run using the standard
single-box moving-refinement mechanism can force-refine only one of its $N$
halos. Everything reported here uses the former.

\subsection{Cost model}
\label{sec:cost}

\todo{Present the six-halo measurement.} Cost separates into an
$N$-independent root grid and a marginal per-halo term:

\begin{deluxetable}{lccc}
\tablecaption{Measured cost decomposition, six halos, DM only
\label{tab:cost}}
\tablehead{\colhead{Level} & \colhead{Root grid} & \colhead{Marginal halo} &
           \colhead{Standalone}\\
           \colhead{} & \colhead{(core-h)} & \colhead{(core-h)} &
           \colhead{(core-h)}}
\startdata
L1 & 17.8 & 8.6  & 55.8 \\
L2 & 18.0 & 15.6 & 73.6 \\
L3 & 19.6 & 70.6 & 129.2 \\
\enddata
\tablecomments{Six halos evolved together cost \measured{0.40$\times$} six
standalone runs.}
\end{deluxetable}

The shared-timestep tax does \emph{not} grow with $N$:
\measured{$1.4\times$} at two halos and \measured{1.20--1.50$\times$} at six,
because each level's timestep is set by its single worst grid regardless of
how many halos contribute candidates.

\paragraph{Design consequence.} Groups should be formed by excluding
outliers, not by restricting mass diversity, and \emph{must not} mix target
refinement levels: a group shares one timestep hierarchy, so the deepest
member would set the clock for all.

\pending{$N=10$ scaling point from the tenpack, which will show where the
timestep tax finally overtakes the occupancy gain.}

% ============================================================================
\section{Validation}
\label{sec:validation}
% ============================================================================

The central question is whether halos evolved together are the same halos.
They cannot be bitwise identical --- a different domain perturbs roundoff from
the first timestep and $N$-body evolution amplifies it --- so the test is
whether \emph{structure} agrees.

\subsection{Contamination control}
\label{sec:contamination}

\todo{Write up the L3 contamination arc.} Coarse particles are defined
relative to the finest species present. Promoting a zoom from L2 to L3
re-classifies former fine particles as coarse, so a region that is clean at
L2 can be contaminated at L3.

We escalated the zoom radius floor from 0 to 200 to
\measured{400}~\kpch{} and traced contaminating particles back to their
Lagrangian origin to distinguish clumped companion subhalos from diffuse
infall. Acceptance threshold: \measured{3.5\%} coarse mass fraction within
\Rvir{}. \check{final per-halo status table from
\texttt{figures\_z0/DECONTAM\_STATUS.md}}

\subsection{Dark matter structure}
\label{sec:dm-audit}

Six halos, L3 DM, 15 snapshots at identical redshifts, compared against their
standalone counterparts.

\begin{figure*}
\plotone{figures/sixpack_audit}
\caption{Assembly histories (left), the multizoom-to-standalone \Mvir{} ratio
against redshift (centre), and $z=0$ substructure counts (right).
Median ratio \measured{0.997} over \measured{77} progenitor pairs resolved by
more than 1000 fine particles, with the central 68\% inside
\measured{4\%}.}
\label{fig:audit}
\end{figure*}

\begin{figure*}
\plotone{figures/sixpack_projections}
\caption{Substructure and concentration. Within a halo the two runs share
field of view and color scale. Satellites appear at the same positions, and
the spherically averaged density profiles are indistinguishable from
$0.03\,\Rvir$ outward.}
\label{fig:projections}
\end{figure*}

\todo{Write the two excursions honestly:}
\begin{itemize}
\item $z \le 0.2$: merger timing. halo24122 differs by
      \measured{$-13.8\%$} in \Mvir{}, but its mass within 300~\kpch{} differs
      by only \measured{$-11.4\%$}, and halo48014 is \measured{$+7.2\%$} in
      \Mvir{} against \measured{$-4.6\%$} within 300~\kpch{} --- mass moving
      between main halo and companion, not appearing or vanishing.
\item $z = 8$--10: progenitor identification, not resolution. halo21246 falls
      to \measured{0.46--0.65} while holding \measured{1600--4900} fine
      particles in both runs; the proto-halo is still several comparable
      clumps and the two runs disagree over the main branch.
\end{itemize}

\subsection{The run-to-run noise floor}
\label{sec:noise-floor}

\pending{Two standalone L3 DM runs repeated at 32 ranks instead of 64 ---
identical ICs, physics, and binary, with only the MPI decomposition changed.
This calibrates how much two runs of the same halo differ under a
perturbation unrelated to multizoom, converting the numbers in
\S\ref{sec:dm-audit} from an upper bound on the multizoom effect into a
measurement of it. Halos 42502 (matched to 0.2\%) and 24122 (the outlier).}

\subsection{Baryons}
\label{sec:gas-audit}

\pending{Sixpack L2 gas at nref7 against the six isolated nref7 runs.
Expected to be a harsher test than dark matter: star formation is stochastic
and the H$_2$ prescription is bistable, so a roundoff-level perturbation near
the ignition threshold can shift ignition timing and move \Mstar{} far more
than the 0.3\% seen in halo mass.}

\subsection{Code regression}
\label{sec:regression}

The multizoom patches must be inert for a single zoom, since the same binary
runs the standalone fleet. With $N=1$, identical nested ICs through the
patched and unpatched binaries give byte-identical grid hierarchies and
\emph{exactly} identical particle masses, with position and velocity
differences at double precision epsilon
(\measured{$\max |\Delta| \sim 10^{-16}$}).
The patched build is therefore safe for single-zoom production but is not
bitwise reproducible against the older one, so the binary must be recorded
when comparing runs.

% ============================================================================
\section{Resolution}
\label{sec:resolution}
% ============================================================================

\todo{Write up the controlled nref7 vs nref9 experiment} at fixed L2 dark
matter, \measured{415} matched snapshot pairs across \measured{4} halos.

\begin{figure}
\plotone{figures/nref7_vs_nref9}
\caption{Stellar mass growth at nref7 and nref9 at fixed dark matter
resolution. Gas reservoirs are identical until $z \sim 5$; stellar masses
diverge by \measured{10--200$\times$}.}
\label{fig:resolution}
\end{figure}

\textbf{This is the dominant systematic in the campaign} and it governs how
the results in \S\ref{sec:results} may be read: an nref7 sample and an nref9
sample are not interchangeable, and any comparison --- including the
multizoom validation --- must be made at matched resolution.

% ============================================================================
\section{First results}
\label{sec:results}
% ============================================================================

\subsection{Star formation histories}
\begin{figure}
\plotone{figures/sfh_fleet_z0}
\caption{Star formation histories of the fleet. Halos that form no stars are
carried as flat dotted lines rather than dropped, since they may ignite at
higher resolution.}
\label{fig:sfh}
\end{figure}
\todo{}

\subsection{Scaling relations}
\begin{figure*}
\plotone{figures/scaling_relations_pooled-targets_z0}
\caption{Stellar-to-halo mass, mass--metallicity, star-forming main sequence,
and baryon budget at $z=0$.}
\label{fig:scaling}
\end{figure*}
\todo{}

\subsection{Quenched fraction}
\label{sec:quenched}

\begin{figure}
\plotone{figures/quenched_fraction_z0}
\caption{Quenched fraction against stellar and halo mass, with Wilson
binomial intervals, compared to ELVES-Field and SAGA background samples.}
\label{fig:quenched}
\end{figure}

\paragraph{Definition.} A galaxy is quenched if
$\mathrm{SFR}(100\,\mathrm{Myr}) = 0$. We deliberately do \emph{not} use
specific star formation rate: dividing by \Mstar{} mislabels massive
star-forming galaxies, and a zero-star-particle halo has no defined sSFR at
all. Halos with no star particles are counted as quenched.

\check{Current sample: 6 of 13 quenched, with a sharp shutoff near
$\Mvir \approx 3\times10^{9}\,\Msun$. Re-derive at freeze --- the sample is
growing.}

\todo{Comparison to ELVES-Field (4 of 15 quenched at
$10^6$--$10^7\,\Msun$) and to the SAGA background analysis, which measures
quenched fractions only above $\sim 4\times10^{8}\,\Msun$ --- a mass the
current fleet does not reach, which is itself worth stating.}

\subsection{Environment is not the driver}
\begin{figure}
\plotone{figures/environment_quenching_z0}
\caption{Quenched status against three environment metrics. The fleet was
selected to be isolated; quenched halos are, if anything, \emph{more}
isolated than star-forming ones.}
\label{fig:environment}
\end{figure}
\todo{State the caveat plainly: quenched and low-mass are the same halos in
this sample, so mass and environment are not independent here.}

% ============================================================================
\section{Discussion}
\label{sec:discussion}
% ============================================================================
\todo{}
\begin{itemize}
\item What multizoom does and does not buy. It is an efficiency method, not a
      physics change; the validation bounds how much structure it perturbs.
\item Where the method breaks down: mixed target levels, and the
      $N$-dependence of the timestep tax at large $N$.
\item Path to the proposal's 250--500 dwarf ambition.
\item Limitations: single feedback prescription, no local radiation, one
      cosmology, resolution systematic of \S\ref{sec:resolution}.
\end{itemize}

% ============================================================================
\section{Summary}
\label{sec:summary}
% ============================================================================
\todo{Numbered findings, mirroring the abstract.}

\begin{acknowledgments}
\todo{NASA ATP grant number, NAS/Aitken allocation, Enzo and yt communities.}
This work used resources of the NASA Advanced Supercomputing (NAS) Division
at Ames Research Center.
\end{acknowledgments}

\software{\enzo{} \citep{Bryan2014},
          \grackle{} \citep{Smith2017},
          \music{} \citep{Hahn2011},
          \ahf{} \citep{Knollmann2009},
          yt \citep{Turk2011},
          astropy \citep{Astropy2013},
          numpy \citep{Harris2020},
          matplotlib \citep{Hunter2007}}

% ============================================================================
\appendix

\section{Performance optimization}
\label{app:perf}
\todo{Summarize the optimization campaign: the subgrid-size floor, the
measured per-level budget, and what was tried and rejected.}

\paragraph{On reporting negative results.} A parameter sweep over the subgrid
size floor produced a suggestive \measured{$-10.7\%$} at the smallest value,
at \measured{2$\sigma$} in a paired per-cycle comparison --- but the four
configurations ran on four different physical nodes, and node-to-node
variation on this machine is of the same order as the effect. We report it as
unresolved rather than as a result. \pending{same-node A/B}

\section{Data management and reproducibility}
\label{app:data}
\todo{The halo registry as the single source of truth; the staged pipeline;
provenance of every run through its build ledger.}

% ============================================================================
\bibliography{refs}
\bibliographystyle{aasjournal}

\end{document}
